\section{Assets and the earnings shock as the state}\label{sec:assets}

The household's own state space is the pair $(a,z)$, assets and the earnings shock, with $z$
following a finite Markov chain and cash on hand $m(a,z)=y(z)+Ra$ in the absence of a means test
on current holdings. The question is then about $g_v(a,z)$, next period's assets.

\subsection{Next period's assets rise with this period's}

\begin{theorem}[Monotone saving in assets]\label{thm:azmono}
Let $a<a'$ both lie in the asset region. Then for every continuation $v$ and every $z$,
$g_v(a,z)\le g_v(a',z)$.
\formal{IncomeFluctuation.policyOf\_mono\_of\_lt}
\end{theorem}

This is Theorem~\ref{thm:topkis} read in the household's state space: cash on hand rises strictly
with assets, so optimal saving does. The development already contained this statement
(\lean{policyOf\_mono}) but with concave slices of the continuation as a hypothesis, and a means
test is precisely what destroys concavity, so the hypothesis had to go. The Topkis argument removes
it, at the cost of requiring the two asset levels to be distinct, which is harmless. As in
Section~\ref{sec:cash}, the continuation cancels and only strict concavity of period utility
survives.

Consumption does not inherit it, for the same reason as before: $c=m(a,z)-g_v(a,z)$ with both
terms rising in $a$, and the witness of Theorem~\ref{thm:cons} transfers directly, since at fixed
$z$ cash on hand is an increasing affine function of assets.

\subsection{When saving in assets does fail}

Theorem~\ref{thm:azmono} has one hypothesis that a means test can remove: that cash on hand rises
with assets. Where the test is assessed on assets carried forward, today's budget is untouched and
the theorem applies verbatim. Where it is assessed on assets currently held, the schedule enters
cash on hand directly, and whether it still rises is a question about the taper rate.

\begin{theorem}[A gentle taper is harmless; a cliff is not]\label{thm:taper}
If $0\le\tau<R$ then $a\mapsto Ra+p(a)$ is strictly increasing under the taper
(\lean{cash\_strictMono\_tapered}), and hence, by Theorem~\ref{thm:topkis}, saving is monotone in
assets (\lean{isOptimalSaving\_mono\_assets}). Under a cliff it is not: if $a\le\bar a<b$ and
$R(b-a)<p$ then cash on hand is strictly smaller at $b$ (\lean{cash\_lt\_cash\_of\_cliff}), and
there are two asset levels, the second larger, at which the household saves strictly less.
\formal{OLG.exists\_saving\_decrease\_in\_assets}
\end{theorem}

The threshold is sharp and it is the taper rate against the gross return. A pension withdrawn more
slowly than assets earn leaves cash on hand rising, and then no amount of damage to the value
function can make saving fall; a pension withdrawn faster makes a household with more assets poorer
in the only sense that matters for today's decision, and saving falls with it. The proof of the
first half is that the positive part is one-Lipschitz, so the pension falls by at most $\tau$ per
unit of assets while gross assets rise by $R$ (\lean{max\_zero\_sub\_le}).

The witness for the second half is again exact, and deliberately built on a \emph{concave}
continuation so that nothing is hidden in it. With a unit gross return, a pension of one withdrawn
above assets of one, and continuation $\log(x+1)$:

\begingroup\small
\begin{center}
\begin{tabular}{@{}lccc@{}}
\toprule
Assets $a$ & cash on hand $Ra+p(a)$ & optimal saving & consumption\\
\midrule
$1$ & $2$ & $1/2$ & $3/2$\\
$3/2$ & $3/2$ & $1/4$ & $5/4$\\
\bottomrule
\end{tabular}
\end{center}
\endgroup

The household with more assets has less to spend, saves less and consumes less. The
non-monotonicity here lives in the budget constraint, not in the shape of the value function, and
that is the whole difference between this case and Section~\ref{sec:cash}.
